\section{Conclusion and Future Work}
\label{sec:conclusion}

This paper studied task-agnostic, parameter-efficient continual learning for sequential code generation, a setting in which a model must adapt to a growing stream of generation, translation, summarization, refinement, and multilingual program-synthesis tasks without forgetting earlier ones and without access to task identity at inference time. We proposed GMM-Router, which trains a dedicated LoRA expert per task on a frozen backbone and routes each query using a diagonal Gaussian Mixture Model fitted on routing embeddings extracted through a fixed orthonormal projection, so that per-task density models remain comparable as the task sequence grows; for executable code generation, a second GMM fitted on runtime error tracebacks, combined with the instruction router through a Jensen--Shannon-divergence-gated rule, recovers samples that fail on every Round-1 attempt.

Across two complementary benchmarks --- an eight-task extension of CodeTask-CL scored by text-overlap metrics and a nine-language executable benchmark scored by unit-test execution --- the experiments in Section~\ref{sec:results} answer the four research questions posed in Section~\ref{sec:rqs}. GMM-Router achieves the highest average score (51.62\%) among all task-agnostic continual learning methods on CodeTask, approaching the task-identity upper bound (54.15\%) with near-zero forgetting (BWT $-1.01\%$, FM $1.01\%$; RQ1, RQ3). On the executable benchmark, it reaches 34.40\% Pass@1 and the highest Pass@5 (47.75\%) among all evaluated methods, exceeding the non-iterative upper bound on Pass@1 (RQ2). Router assignment analysis shows a sharp contrast between the two benchmarks: instruction-based hard routing reaches 90.6\% accuracy on the type-diverse CodeTask sequence but only 15.8\% on the syntactically homogeneous executable benchmark, where the complementary traceback router recovers 96.8\% accuracy in Round-2 (RQ3). The iterative refinement pipeline contributes 1.36pp Pass@1 and 6.09 points Pass@5 over Round-1 alone, concentrated on Rust and Java, and JSD gating specifically recovers Java and Shell beyond what traceback routing alone achieves (RQ4). Taken together, these results indicate that per-task expert isolation and density-based probabilistic routing are complementary: isolation prevents forgetting on structurally diverse task sequences, and soft, distribution-level routing with iterative refinement provides candidate diversity and error recovery on structurally homogeneous ones, without requiring task labels, data replay, or a shared adapter subspace.

GMM-Router also has limitations, discussed in Section~\ref{sec:limitations}, that motivate three directions for future work. First, per-task expert isolation limits knowledge transfer between related tasks, most visibly on C\# and Rust in the executable benchmark, where shared-adaptation baselines benefit from cross-language transfer that GMM-Router cannot currently exploit; a promising direction is a lightweight, jointly trained module shared beneath the task-specific adapter layers, or a distillation step between adapters for tasks whose GMMs the router already identifies as overlapping, turning the density model into a signal for both routing and selective knowledge consolidation. Second, adapter storage grows linearly with the number of tasks; dynamic expert merging --- pruning or merging adapters whose GMMs exhibit high overlap in the routing space --- or post-hoc adapter compression could keep GMM-Router practical for task sequences much longer than the eight- and nine-task streams studied here. Third, the repair pairs used for adapter refinement in Section~\ref{sec:refinement} do not guarantee that the target solution resolves the specific error shown in the accompanying traceback, which is a plausible cause of the Rust regression observed when refinement is combined with JSD-gated routing (Table~\ref{tab:ablation_refinement}); constructing repair pairs with execution-guided search for patches that provably fix the reported error is a natural next step and may recover Rust's gains while preserving the improvement already observed on Java.
